\section{Method}

\subsection{Overview}

We propose \textbf{DreamerVLA}, a Dreamer-style training framework that combines the semantic prior of a pretrained Vision-Language-Action (VLA) model with the long-horizon optimization ability of latent world models. The current DreamerVLA pipeline uses a pi0-style action-query VLA head to produce action-conditioned hidden observations, trains a DreamerV3 recurrent state-space model (RSSM) over these hidden observations, and then optimizes an actor-critic policy through imagined latent rollouts.

At each time step $t$, the robot observes two RGB views, a language instruction, and proprioceptive state:
\[
o_t = \{I_t^{\text{third}}, I_t^{\text{wrist}}, x_t^{\text{prop}}, l\}.
\]
The VLA model first converts this multimodal observation into an action-query hidden sequence
\[
u_t = f_{\text{VLA}}^{\text{query}}(o_t) \in \mathbb{R}^{K \times 1024},
\]
where $K$ is the action horizon. We flatten this sequence into the world-model observation
\[
e_t = \mathrm{vec}(u_t) \in \mathbb{R}^{5120},
\]
using $K=5$ in our LIBERO configuration. These hidden observations are precomputed and stored in sidecar HDF5 files aligned with the original demonstrations, so world-model and actor-critic training do not backpropagate through the large VLA backbone.

DreamerVLA contains three trainable components:
\begin{enumerate}
    \item an \textbf{action-hidden DreamerV3 world model} that uses pi0 action-query hidden states as observations for a discrete RSSM while reconstructing pixels and action-hidden observations;
    \item a \textbf{pi0 action-hidden actor} that maps imagined RSSM features back to action-query hidden states and reuses the pretrained pi0 output projection to produce action chunks;
    \item a \textbf{two-hot critic} trained from DreamerV3-style imagined $\lambda$-returns.
\end{enumerate}

This design keeps the VLA backbone fixed as a semantic action-conditioned encoder, learns compact task dynamics in latent space, and optimizes the action head through imagined rollouts without additional environment interaction.

\subsection{Overall Training Pipeline}

DreamerVLA is trained in four stages. First, we prepare LIBERO demonstrations by removing no-op transitions and storing synchronized third-person RGB, wrist RGB, robot states, actions, rewards, and terminal labels. The action horizon is set to $K=5$, and each action is represented as a 7-dimensional continuous control vector.

Second, we warm-start the VLA policy with supervised action prediction. Starting from the pretrained RynnVLA/Chameleon backbone, we fine-tune a pi0-style action-query head on LIBERO demonstrations. Instead of predicting actions directly from all language-image tokens, the head appends one learned action query for each future action step and regresses an action chunk from the resulting query hidden states. This stage provides the supervised VLA checkpoint used by all downstream DreamerVLA components.

Third, we freeze the VLA backbone and precompute action-conditioned hidden observations for every frame. The resulting sidecar files store $e_t$ together with metadata for the VLA checkpoint, action horizon, action-head type, and hidden-source type. The dataset loader checks this metadata at startup, preventing stale hidden sidecars from being mixed with incompatible checkpoints.

Fourth, we train the action-hidden DreamerV3 world model and then optimize the actor and critic with imagined rollouts. World-model pretraining can be run as a standalone stage, and the resulting checkpoint can be loaded into joint DreamerVLA training. In the actor-critic stage, posterior latent states inferred from replay windows initialize imagined trajectories; the actor produces action chunks from imagined hidden states, and the first action in each chunk advances the RSSM rollout.

\subsection{Frozen VLA Action-Query Features}

We use the pretrained RynnVLA/Chameleon model as a frozen multimodal encoder. For each frame, the third-person view, wrist-camera view, task text, and optional proprioceptive state are serialized into the prompt format used by the VLA model. The Chameleon image tokenizer and text tokenizer convert the input into a unified token sequence, which is processed by the frozen Transformer backbone.

The pi0-style action-query head appends $K$ learned query tokens, one per future action step. Context tokens condition these queries through attention, and the head exposes the query outputs as compact action-conditioned hidden states:
\[
u_t = [u_t^1,\ldots,u_t^K],
\qquad
u_t^k \in \mathbb{R}^{1024}.
\]
The world model uses their flattened representation $e_t=\mathrm{vec}(u_t)$ as its observation. Compared with storing a full VLA token sequence, this representation is smaller and directly aligned with the downstream action head. During Dreamer training, the VLA backbone remains frozen and is not included in the optimization graph.

\subsection{Action-Hidden DreamerV3 World Model}

The world model follows the DreamerV3 recurrent state-space model, but replaces the standard pixel encoder with a lightweight encoder over precomputed action-hidden observations. Given action-hidden observations $\{e_t\}_{t=1}^{T}$ and actions $\{a_t\}_{t=1}^{T}$, the encoder maps each $e_t$ into the RSSM observation space. The RSSM state is
\[
s_t = (h_t, z_t),
\]
where $h_t \in \mathbb{R}^{8192}$ is a deterministic recurrent state and $z_t$ is a discrete stochastic state with $32$ variables and $64$ categories per variable. The flattened latent feature is
\[
\phi_t = [h_t; \mathrm{vec}(z_t)] \in \mathbb{R}^{10240}.
\]

\paragraph{RSSM transition and posterior.}
The deterministic transition is a DreamerV3 block-GRU update conditioned on the previous deterministic state, stochastic state, and action:
\[
h_t = g_\theta(h_{t-1}, z_{t-1}, a_t).
\]
The model predicts a prior distribution $p_\theta(z_t \mid h_t)$ and an observation-conditioned posterior
\[
q_\theta(z_t \mid h_t, e_t).
\]
We use the DreamerV3 unimix categorical parameterization and split the KL regularizer into dynamics and representation terms:
\[
\mathcal{L}_{\text{dyn}} =
\mathrm{KL}\big(\mathrm{sg}(q_\theta(z_t \mid h_t,e_t)) \,\|\, p_\theta(z_t \mid h_t)\big),
\]
\[
\mathcal{L}_{\text{rep}} =
\mathrm{KL}\big(q_\theta(z_t \mid h_t,e_t) \,\|\, \mathrm{sg}(p_\theta(z_t \mid h_t))\big).
\]

\paragraph{Multi-target reconstruction.}
The RSSM is trained with both visual and action-hidden reconstruction losses. A pixel decoder reconstructs the concatenated third-person and wrist images from $\phi_t$:
\[
\hat{I}_t = d_{\text{img}}(\phi_t),
\qquad
\mathcal{L}_{\text{img}} = \|\hat{I}_t - I_t\|_2^2.
\]
A hidden decoder predicts the flattened pi0 action-query hidden observation:
\[
\hat{e}_t = d_{\text{hid}}(\phi_t),
\qquad
\mathcal{L}_{\text{hid}} = \|\hat{e}_t - e_t\|_2^2.
\]
This hidden reconstruction is the bridge between the RSSM and the pi0 action-hidden actor: during imagination, the actor consumes $\hat{e}_t$ rather than raw pixels or full VLA token states.

\paragraph{Reward and continuation.}
The world model predicts sparse task reward and continuation from the latent feature:
\[
\hat{r}_t = d_r(\phi_t), \qquad \hat{c}_t = d_c(\phi_t).
\]
For LIBERO-style sparse success rewards, the implementation uses a binary reward head. The continuation head is trained with a binary cross-entropy objective over non-terminal targets.

The full world-model objective is
\[
\mathcal{L}_{\text{WM}}
=
\lambda_{\text{img}}\mathcal{L}_{\text{img}}
+ \lambda_{\text{dyn}}\mathcal{L}_{\text{dyn}}
+ \lambda_{\text{rep}}\mathcal{L}_{\text{rep}}
+ \lambda_{\text{hid}}\mathcal{L}_{\text{hid}}
+ \lambda_r\mathcal{L}_{r}
+ \lambda_c\mathcal{L}_{c}.
\]

\subsection{Pi0 Action-Hidden Actor}

Instead of using a small MLP policy, DreamerVLA reuses the action prior learned by the supervised pi0-style VLA head. The world model first decodes an imagined latent feature into a reconstructed action-hidden vector,
\[
\hat{e}_t = d_{\text{hid}}(\phi_t) \in \mathbb{R}^{5120},
\]
which is reshaped into
\[
\hat{u}_t \in \mathbb{R}^{K \times 1024}.
\]
A lightweight adapter refines this sequence, and the pretrained pi0 output projection maps each action-query hidden state to a continuous action. The actor therefore outputs an action chunk
\[
\mu_{t:t+K-1} = \pi_\psi(\hat{u}_t) \in \mathbb{R}^{K \times A}.
\]
In our LIBERO configuration, $K=5$ and $A=7$. During imagination, the first action in the chunk is used as the action mean,
\[
\mu_t = \mu_{t:t+K-1}[0],
\]
and a learned diagonal log-standard deviation defines a Gaussian policy:
\[
a_t \sim \mathcal{N}(\mu_t, \mathrm{diag}(\sigma^2)).
\]
The actor can freeze or fine-tune the pi0 output projection. In both cases, the policy is initialized from the supervised VLA checkpoint and further optimized by DreamerV3 actor-critic updates.

\subsection{Imagination-Based Actor-Critic Learning}

Actor-critic training starts from posterior latent states inferred from replay sequences. From each start state, the world model generates an imagined trajectory of horizon $H$:
\[
a_t \sim \pi_\psi(\cdot \mid \phi_t),
\qquad
s_{t+1} \sim p_\theta(\cdot \mid s_t,a_t),
\qquad
\hat{r}_{t+1}=d_r(\phi_{t+1}), \quad
\hat{c}_{t+1}=d_c(\phi_{t+1}).
\]
Gradients are stopped through imagined world-model states, matching the standard DreamerV3 actor update.

The critic receives the RSSM feature $\phi_t$ and predicts a categorical distribution over symlog value bins. Its scalar value is the symexp-transformed expectation:
\[
V_\omega(\phi_t)=\mathrm{symexp}\left(\mathbb{E}_{b\sim p_\omega(\cdot\mid \phi_t)}[b]\right).
\]
A Polyak-averaged target critic provides bootstrap values. We compute continuation-weighted $\lambda$-returns:
\[
R_t^\lambda =
\hat{r}_{t+1}
+ \gamma \hat{c}_{t+1}
\left((1-\lambda)V_{\bar{\omega}}(\phi_{t+1})
+ \lambda R_{t+1}^{\lambda}\right).
\]

The actor is trained with a likelihood-ratio objective using the imagined advantage:
\[
\mathcal{L}_{\pi}
=
-
\mathbb{E}\left[
w_t
\left(
\log \pi_\psi(a_t\mid \phi_t)\,
\mathrm{sg}\left(\frac{R_t^\lambda - V_\omega(\phi_t)}{S}\right)
+ \eta \mathcal{H}[\pi_\psi(\cdot\mid \phi_t)]
\right)
\right],
\]
where $w_t$ is the cumulative continuation discount and $S=\max(1,P_{95}^{\text{ema}}-P_{5}^{\text{ema}})$ is the running percentile return scale. The critic is trained by maximizing the log probability of the two-hot symlog target:
\[
\mathcal{L}_{V}
=
-
\mathbb{E}\left[
\log p_\omega\left(\mathrm{twohot}(\mathrm{sg}(R_t^\lambda)) \mid \phi_t\right)
\right].
\]

\subsection{Training Strategy}

Training alternates between two phases on the same replay windows. First, the world model is updated using real image sequences, precomputed action-hidden observations, actions, rewards, and terminal labels. Second, the actor and critic are updated from imagined rollouts produced by the learned RSSM. Only the world model, pi0 action-hidden actor, and critic are optimized; the VLA backbone used to generate action-query hidden observations remains frozen throughout.
